% =============================================================================
% Capacitação — Git, GitLab e IA aplicada à Engenharia de Software
% Apresentação expositiva (1 h). Compilar com:  latexmk -xelatex apresentacao.tex
% =============================================================================
\documentclass[aspectratio=169,11pt]{beamer}
% =============================================================================
% marinha-comum.tex — identidade visual compartilhada pela apresentação e pela
% apostila da capacitação. Segue o Manual de Identidade Visual da Marinha do
% Brasil (dez/2025): cores institucionais (p. 23) e tipografia (pp. 24-30).
%
% Compilar sempre com XeLaTeX (fontspec). Os caminhos são relativos à pasta
% do documento (apresentacao/ ou apostila/).
% =============================================================================

% ---------- Caminhos ---------------------------------------------------------
\newcommand{\raizcapacitacao}{..}
\newcommand{\imagens}{\raizcapacitacao/assets/imagens}
\newcommand{\fontes}{\raizcapacitacao/assets/fontes}

% ---------- Cores institucionais (Manual, p. 23) -----------------------------
\usepackage{xcolor}
\usepackage{graphicx}
\definecolor{azulMB}{HTML}{050F41}   % Azul Marinha  · Pantone 2768C
\definecolor{ouroMB}{HTML}{FAB932}   % Amarelo Ouro  · Pantone 1235C
\definecolor{verdeMB}{HTML}{079551}  % Verde Brasil  · Pantone 355C
% Tons de apoio derivados das institucionais (fundos de caixas e diagramas)
\colorlet{azulClaro}{azulMB!8}
\colorlet{azulMedio}{azulMB!55}
\colorlet{ouroClaro}{ouroMB!18}
\colorlet{verdeClaro}{verdeMB!12}
\definecolor{cinzaTexto}{HTML}{3A3F55}
\definecolor{cinzaSuave}{HTML}{F3F4F8}
\definecolor{vermelhoAlerta}{HTML}{B3261E}
\colorlet{vermelhoClaro}{vermelhoAlerta!10}

% ---------- Tipografia (Manual, pp. 24-30) -----------------------------------
% Principal: Octin College (títulos) e alternativa Bebas Neue Pro são fontes
% comerciais. Usamos a Bebas Neue livre (OFL), da mesma família de desenho,
% nos títulos, e a Montserrat (tipografia de apoio oficial, OFL) no texto.
\usepackage{fontspec}
\setmainfont{Montserrat}[
  Path=\fontes/montserrat/, Extension=.otf,
  UprightFont=*-Regular, BoldFont=*-SemiBold,
  ItalicFont=*-Italic, BoldItalicFont=*-SemiBoldItalic]
\setsansfont{Montserrat}[
  Path=\fontes/montserrat/, Extension=.otf,
  UprightFont=*-Regular, BoldFont=*-SemiBold,
  ItalicFont=*-Italic, BoldItalicFont=*-SemiBoldItalic]
\setmonofont{FiraMono}[
  Path=\fontes/fira-mono/, Extension=.otf,
  UprightFont=*-Regular, BoldFont=*-Bold, Scale=0.9]
\newfontfamily\titulofonte{BebasNeue-Regular}[
  Path=\fontes/bebas-neue/, Extension=.ttf, LetterSpace=2.0]
\newfontfamily\montblack{Montserrat-Black}[Path=\fontes/montserrat/, Extension=.otf]
\newfontfamily\montlight{Montserrat-Light}[Path=\fontes/montserrat/, Extension=.otf]

% ---------- Pacotes comuns ----------------------------------------------------
\usepackage{tikz}
\usetikzlibrary{positioning,arrows.meta,calc,shapes.geometric,fit,backgrounds,shadows.blur,decorations.pathmorphing}
\usepackage{booktabs,tabularx,array}
\usepackage{fontawesome5}
\usepackage{listings}
\usepackage{csquotes}

% ---------- Código (listings) -------------------------------------------------
\lstdefinestyle{terminal}{
  basicstyle=\ttfamily\small\color{cinzaTexto},
  columns=fixed, basewidth=0.6em, keepspaces=true,
  breaklines=true, breakatwhitespace=true,
  showstringspaces=false, upquote=true,
  commentstyle=\color{verdeMB},
  keywordstyle=\color{azulMB}\bfseries,
  morekeywords={git,glab,claude},
  morecomment=[l]{\#},
  literate=
    {á}{{\'a}}1 {à}{{\`a}}1 {ã}{{\~a}}1 {â}{{\^a}}1
    {é}{{\'e}}1 {ê}{{\^e}}1 {í}{{\'i}}1 {ó}{{\'o}}1
    {õ}{{\~o}}1 {ô}{{\^o}}1 {ú}{{\'u}}1 {ç}{{\c c}}1
    {Á}{{\'A}}1 {É}{{\'E}}1 {Í}{{\'I}}1 {Ó}{{\'O}}1
    {Ú}{{\'U}}1 {Ç}{{\c C}}1 {Ã}{{\~A}}1 {Õ}{{\~O}}1
    {→}{{$\rightarrow$}}1 {—}{{---}}1 {·}{{$\cdot$}}1
}
\lstdefinestyle{saida}{style=terminal, keywordstyle=, commentstyle=,
  basicstyle=\ttfamily\small\color{cinzaTexto}}

% ---------- Diagramas de histórico Git (TikZ) ---------------------------------
\tikzset{
  commit/.style={circle, draw=#1, fill=white, line width=1.6pt,
    minimum size=7mm, inner sep=0pt, font=\small\bfseries\color{#1}},
  commit/.default=azulMB,
  commitm/.style={circle, draw=ouroMB!80!black, fill=ouroMB, line width=1.6pt,
    minimum size=8mm, inner sep=0pt, font=\small\bfseries\color{azulMB}},
  hist/.style={line width=1.8pt, color=#1},
  hist/.default=azulMB,
  rotulo/.style={rounded corners=2pt, fill=#1, text=white,
    font=\scriptsize\bfseries, inner sep=3pt},
  rotulo/.default=azulMB,
  caixa/.style={rounded corners=3pt, draw=#1, line width=1pt, fill=#1!8,
    align=center, font=\small, inner sep=5pt, text=azulMB},
  caixa/.default=azulMB,
  seta/.style={-{Stealth[length=3mm]}, line width=1.2pt, color=azulMB},
  setafina/.style={-{Stealth[length=2mm]}, line width=0.8pt, color=azulMedio},
}

% ---------- Marca do projeto (texto) -------------------------------------------
% Usada onde a versão interna tinha o logotipo da Marinha, que não entra na
% versão pública. \marca[cor das letras]{tamanho em pt}
\newcommand{\marca}[2][azulMB]{{\titulofonte\fontsize{#2}{#2}\selectfont
  \textcolor{#1}{GIT}\textcolor{ouroMB}{\,·\,}\textcolor{#1}{IA}\textcolor{ouroMB}{\,·\,}\textcolor{#1}{DOCS}}}

\usepackage{polyglossia}
\setdefaultlanguage[variant=brazilian]{portuguese}
\usepackage{pgfplots}
\pgfplotsset{compat=1.18}

% ---------- Tema (Manual de Identidade Visual da MB) -------------------------
\usefonttheme{professionalfonts}
\setbeamertemplate{navigation symbols}{}
\setbeamercolor{normal text}{fg=cinzaTexto,bg=white}
\setbeamercolor{structure}{fg=azulMB}
\setbeamercolor{alerted text}{fg=verdeMB}
\setbeamerfont{alerted text}{series=\bfseries}
\setbeamertemplate{itemize item}{\color{ouroMB}\raisebox{0.15ex}{\scriptsize$\blacksquare$}}
\setbeamertemplate{itemize subitem}{\color{verdeMB}\scriptsize$\blacktriangleright$}
\setbeamertemplate{enumerate item}{\color{azulMB}\bfseries\insertenumlabel.}
\setlength{\leftmargini}{1.4em}

% Faixa superior azul com o título em Bebas Neue e o logotipo branco à direita
\setbeamertemplate{frametitle}{%
  \nointerlineskip
  \begin{tikzpicture}[remember picture, overlay]
    \fill[azulMB] (current page.north west) rectangle ([yshift=-11mm]current page.north east);
    \fill[ouroMB] ([yshift=-11mm]current page.north west) rectangle ([yshift=-11.8mm]current page.north east);
    \node[anchor=west, text=white, font=\titulofonte\fontsize{21}{21}\selectfont]
      at ([xshift=8mm,yshift=-5.7mm]current page.north west) {\MakeUppercase{\insertframetitle}};
    \node[anchor=east, inner sep=0pt]
      at ([xshift=-7mm,yshift=-5.5mm]current page.north east)
      {\marca[white]{15}};
  \end{tikzpicture}%
  \vspace*{9mm}%
}
\setbeamertemplate{footline}{%
  \begin{tikzpicture}[remember picture, overlay]
    \node[anchor=south west, font=\tiny\color{azulMedio}]
      at ([xshift=8mm,yshift=2.5mm]current page.south west)
      {Git, GitLab e IA aplicada à Engenharia de Software \,·\, Marinha do Brasil \,·\, 01/10/2026};
    \node[anchor=south east, font=\tiny\bfseries\color{azulMB}]
      at ([xshift=-7mm,yshift=2.5mm]current page.south east)
      {\insertframenumber\,/\,\inserttotalframenumber};
  \end{tikzpicture}%
}
\setbeamersize{text margin left=8mm, text margin right=8mm}

% Divisória de seção: fundo azul, número grande em ouro (abertura de capítulo
% no mesmo espírito do manual, p. 31)
\newcommand{\secao}[3]{%
  \section{#2}
  {\setbeamertemplate{footline}{}
   \begin{frame}[plain]
   \begin{tikzpicture}[remember picture, overlay]
     \fill[azulMB] (current page.south west) rectangle (current page.north east);
     \fill[ouroMB] ([xshift=14mm,yshift=-3mm]current page.west) rectangle ++(28mm,1.2mm);
     \node[anchor=south west, text=ouroMB, font=\titulofonte\fontsize{96}{96}\selectfont]
       at ([xshift=13mm,yshift=2mm]current page.west) {#1};
     \node[anchor=north west, text=white, font=\titulofonte\fontsize{40}{40}\selectfont, align=left]
       at ([xshift=14mm,yshift=-8mm]current page.west) {\MakeUppercase{#2}};
     \node[anchor=north west, text=white!80!azulMB, font=\montlight\large, align=left, text width=11cm]
       at ([xshift=14mm,yshift=-26mm]current page.west) {#3};
     \node[anchor=south east, inner sep=0pt]
       at ([xshift=-8mm,yshift=7mm]current page.south east)
       {\marca[white]{20}};
   \end{tikzpicture}
   \end{frame}}%
}

% Blocos
\newcommand{\destaque}[2][ouroMB]{%
  \begin{tikzpicture}\node[fill=#1!18, draw=#1, line width=0.8pt, rounded corners=3pt,
    inner sep=7pt, text width=\linewidth-16pt, font=\small, text=azulMB] {#2};\end{tikzpicture}}
\newcommand{\etiqueta}[2][azulMB]{\tikz[baseline=(e.base)]\node[rotulo=#1](e){#2};}
\newcommand{\cmd}[1]{{\ttfamily\color{azulMB}#1}}
\lstset{style=terminal}

\title{Git, GitLab e IA aplicada à Engenharia de Software}
\author{Ronivaldo D. Andrade}
\date{01/10/2026}

\begin{document}

% =============================================================================
% CAPA
% =============================================================================
{\setbeamertemplate{footline}{}
\begin{frame}[plain]
\begin{tikzpicture}[remember picture, overlay]
  % A foto (16:9) é mais larga que o painel: desloca para a ilha ficar à vista.
  \node[anchor=north east, inner sep=0pt] at ([xshift=5.2cm]current page.north east)
    {\includegraphics[height=\paperheight]{\imagens/ilha_das_cobras.jpg}};
  \fill[azulMB] (current page.south west) rectangle ([xshift=0.6\paperwidth]current page.north west);
  \fill[ouroMB] ([xshift=0.6\paperwidth]current page.south west) rectangle ([xshift=0.605\paperwidth]current page.north west);
  \node[anchor=north west, inner sep=0pt] at ([xshift=10mm,yshift=-8mm]current page.north west)
    {\marca[white]{20}};
  \node[anchor=north west, text=ouroMB, font=\small\bfseries] (k) at ([xshift=10mm,yshift=-22mm]current page.north west)
    {CAPACITAÇÃO INTERNA};
  \node[anchor=north west, text=white, align=left, inner xsep=0pt,
        font=\titulofonte\fontsize{25}{26}\selectfont] (t) at ([yshift=-1mm]k.south west)
    {GIT, GITLAB E IA APLICADA\\À ENGENHARIA DE SOFTWARE};
  \node[anchor=north west, text=white!80!azulMB, align=left, text width=0.5\paperwidth, inner xsep=0pt,
        font=\montlight\footnotesize] at ([yshift=-2mm]t.south west)
    {Versionamento, colaboração e documentação de sistemas legados a partir de dois problemas reais};
  \node[anchor=south west, text=white, align=left, font=\small] at ([xshift=10mm,yshift=6mm]current page.south west)
    {\textbf{Ronivaldo D. Andrade}\\[1pt]
     {\color{white!75!azulMB}\scriptsize Estagiário de Análise e Desenvolvimento de Sistemas}\\
     {\color{white!75!azulMB}\scriptsize Marinha do Brasil}\\[2pt]
     {\color{ouroMB}\footnotesize\bfseries 01/10/2026 \,·\, 10h às 11h}};
\end{tikzpicture}
\end{frame}}

% =============================================================================
\begin{frame}{O que vamos ver nesta hora}
\begin{columns}[T]
\column{0.5\textwidth}
\begin{tikzpicture}[every node/.style={anchor=west}]
  \foreach \n/\t/\s/\y in {
      01/Fundamentos/{Git, GitLab, branch, Merge Request, glab}/0,
      02/Problema 1/{Trazer um sistema já versionado para o GitLab institucional}/-1.55,
      03/Problema 2/{Documentar sistemas legados com um agente de IA}/-3.1,
      04/Encerramento/{A visão geral e o que levar daqui}/-4.65}{
    \node[font=\titulofonte\fontsize{30}{30}\selectfont, text=ouroMB] at (0,\y) {\n};
    \node[font=\bfseries\color{azulMB}] at (1.25,\y+0.22) {\t};
    \node[font=\footnotesize, text width=5.4cm] at (1.25,\y-0.32) {\s};
  }
\end{tikzpicture}
\column{0.46\textwidth}
\vspace{1mm}
{\small\bfseries\color{azulMB} Ao final, você deve conseguir:}
\begin{itemize}\small
  \item diferenciar \alert{Git}, \alert{GitLab}, \alert{GitHub} e \alert{glab};
  \item explicar branch, commit, merge e Merge Request;
  \item entender e resolver um conflito;
  \item unir dois históricos Git independentes;
  \item usar um agente de IA como ferramenta de engenharia, com validação humana.
\end{itemize}
\vspace{2mm}
\destaque{\faBook\enspace Os exercícios práticos estão na \textbf{apostila}. Aqui é só a exposição.}
\end{columns}
\end{frame}

% =============================================================================
\secao{01}{Fundamentos}{Git, GitLab, branches, Merge Requests e o GitLab CLI}
% =============================================================================

\begin{frame}{Por que versionar?}
\begin{columns}[c]
\column{0.47\textwidth}
{\small Sem controle de versão, o histórico vira isto:}\\[3mm]
\begin{tikzpicture}[doc/.style={draw=azulMedio, fill=cinzaSuave, rounded corners=2pt,
    font=\ttfamily\scriptsize, anchor=west, minimum width=5.6cm, inner sep=4pt}]
  \node[doc] (a) at (0,0) {\faFile*[regular]\; sistema.zip};
  \node[doc] (b) at (0,-0.62) {\faFile*[regular]\; sistema\_v2.zip};
  \node[doc] (c) at (0,-1.24) {\faFile*[regular]\; sistema\_v2\_final.zip};
  \node[doc] (d) at (0,-1.86) {\faFile*[regular]\; sistema\_v2\_final\_AGORA\_VAI.zip};
  \node[doc, draw=vermelhoAlerta, fill=vermelhoClaro] (e) at (0,-2.48) {\faFile*[regular]\; sistema\_v2\_final\_corrigido(1).zip};
\end{tikzpicture}\\[2mm]
{\footnotesize\color{vermelhoAlerta}\faTimesCircle\; Quem mudou? O quê? Por quê? Qual é a certa?}
\column{0.49\textwidth}
{\small Com o Git, cada mudança fica registrada com:}
\begin{itemize}\small
  \item \textbf{o que} mudou (diferença linha a linha);
  \item \textbf{quem} mudou e \textbf{quando};
  \item \textbf{por que} mudou (mensagem do commit);
  \item a possibilidade de \textbf{voltar} a qualquer estado anterior;
  \item várias pessoas trabalhando \textbf{em paralelo}.
\end{itemize}
\end{columns}
\end{frame}

\begin{frame}[fragile]{O que é o Git?}
\begin{columns}[c]
\column{0.55\textwidth}
\includegraphics[height=9mm]{\imagens/git-logo.pdf}\\[3mm]
{\large Um \alert{sistema distribuído de controle de versão}.}\\[3mm]
\begin{itemize}\small
  \item Cada \textbf{commit} é uma fotografia do projeto num momento.
  \item Os commits formam uma \textbf{linha do tempo}: o histórico.
  \item Funciona \textbf{100\,\% local}, sem servidor nenhum:
\end{itemize}
\vspace{-1mm}
\begin{lstlisting}
git init
git add .
git commit -m "Primeiro commit"
\end{lstlisting}
\column{0.42\textwidth}
\centering
\begin{tikzpicture}
  \foreach \l/\x in {A/0,B/1.3,C/2.6,D/3.9}{\node[commit] (\l) at (\x,0) {\l};}
  \draw[hist] (A) -- (B) -- (C) -- (D);
  \node[rotulo=verdeMB, above=4mm of D] (h) {main};
  \draw[setafina, verdeMB] (h) -- (D);
  \node[font=\scriptsize, text width=5cm, align=center, below=5mm of B.south east] {cada círculo = um commit\\(quem, quando, o quê, por quê)};
\end{tikzpicture}
\end{columns}
\end{frame}

\begin{frame}{Um aviso honesto antes de continuar}
\begin{columns}[c]
\column{0.36\textwidth}
\centering
\includegraphics[height=0.78\textheight]{\imagens/xkcd-1597-git.png}
\column{0.58\textwidth}
{\small\itshape\color{azulMedio}
\enquote{Isto é o Git. Ele controla trabalho colaborativo em projetos por meio de um belo modelo distribuído de árvore da teoria dos grafos.}\\[1mm]
\enquote{Legal. Como a gente usa?}\\[1mm]
\enquote{Não faço ideia. Só decora estes comandos\ldots{} se der erro, salva teu trabalho em outro lugar, apaga o projeto e baixa uma cópia nova.}}\\[5mm]
\destaque[verdeMB]{Hoje a proposta é o contrário: \textbf{entender o modelo} por trás dos comandos, para não precisar apagar nada.}
\vfill
{\tiny\color{azulMedio} xkcd \#1597 “Git”, de Randall Munroe — CC BY-NC 2.5 — tradução livre.}
\end{columns}
\end{frame}

\begin{frame}{Git, GitLab, GitHub e glab não são a mesma coisa}
\renewcommand{\arraystretch}{1.25}
\footnotesize
\begin{tabularx}{\linewidth}{@{}m{14mm} >{\bfseries\color{azulMB}}m{24mm} >{\raggedright\arraybackslash}X >{\raggedright\arraybackslash}X@{}}
\toprule
 & & \textbf{O que é} & \textbf{Para que serve} \\
\midrule
\includegraphics[height=6mm]{\imagens/git-logo.pdf} & Git & Sistema distribuído de controle de versão & Controlar o histórico do código \\
\includegraphics[height=7mm]{\imagens/gitlab-tanuki.pdf} & GitLab & Plataforma baseada em Git, \alert{instalável na própria infraestrutura} & Hospedar, revisar, integrar (CI/CD) e colaborar \\
\includegraphics[height=6mm]{\imagens/github-mark.pdf} & GitHub & Plataforma baseada em Git, majoritariamente em nuvem & Hospedar, revisar e colaborar \\
\centering\faTerminal & GitLab CLI \texttt{glab} & Ferramenta de linha de comando & Operar o GitLab pelo terminal \\
\bottomrule
\end{tabularx}\\[2mm]
\destaque{\textbf{Git} é o motor. \textbf{GitLab} e \textbf{GitHub} são a garagem com oficina, vistoria e controle de acesso. \textbf{glab} é o controle remoto da garagem.}
\end{frame}

\begin{frame}{GitLab: uma plataforma inteira em volta do Git}
\centering
\begin{tikzpicture}[f/.style={caixa, minimum width=3.1cm, minimum height=1.25cm, font=\footnotesize}]
  \foreach \i/\ic/\t in {0/\faDatabase/Repositórios, 1/\faCodeBranch/Branches, 2/\faCodeBranch/Merge Requests, 3/\faSearch/Revisão de código,
                         4/\faUserShield/Permissões, 5/\faTasks/Issues, 6/\faCogs/CI/CD e pipelines, 7/\faTag/Releases}{
    \pgfmathtruncatemacro{\c}{mod(\i,4)}\pgfmathtruncatemacro{\r}{\i/4}
    \node[f] at (\c*3.45, -\r*1.55) {{\color{verdeMB}\large\ic}\\[1pt] \t};
  }
  \node[f, draw=ouroMB, fill=ouroClaro, minimum width=13.45cm] at (5.175,-3.1) {{\color{azulMB}\faClipboardCheck}\; Auditoria, segurança e colaboração entre equipes};
\end{tikzpicture}
\end{frame}

\begin{frame}{GitLab Self-Managed: os dados ficam em casa}
\begin{columns}[T]
\column{0.5\textwidth}
\begin{tikzpicture}
  \node[draw=azulMB, line width=1.2pt, rounded corners=6pt, minimum width=6.4cm, minimum height=4.6cm,
        fill=azulClaro] (om) at (0,0) {};
  \node[rotulo, anchor=north west] at ([xshift=2mm,yshift=-2mm]om.north west) {\faBuilding\; Infraestrutura da instituição};
  \node[caixa=verdeMB, minimum width=2.4cm] (srv) at (0,0.2) {\includegraphics[height=6mm]{\imagens/gitlab-tanuki.pdf}\\GitLab};
  \node[caixa, minimum width=1.6cm, font=\scriptsize] (u1) at (-2,-1.5) {\faUser\\equipe};
  \node[caixa, minimum width=1.6cm, font=\scriptsize] (u2) at (2,-1.5) {\faServer\\backups};
  \draw[setafina,<->] (u1) -- (srv); \draw[setafina,<->] (u2) -- (srv);
  \node[font=\scriptsize\color{vermelhoAlerta}, align=center] at (4.1,1.4) {\faCloud\\internet\\externa};
  \draw[vermelhoAlerta, line width=1.2pt, dashed] (3.15,0.2) -- (3.15,2.2);
\end{tikzpicture}
\column{0.46\textwidth}
{\small A instituição instala e administra \textbf{a própria instância}. Ganha controle sobre:}
\begin{itemize}\footnotesize
  \item código-fonte e histórico;
  \item usuários, permissões e autenticação;
  \item auditoria e políticas de segurança.
\end{itemize}
\destaque[vermelhoAlerta]{\faExclamationTriangle\; \textbf{Estar dentro de casa não é estar seguro.} Segurança continua dependendo de configuração, atualização, backup, firewall, monitoramento e controle de acesso.}
\end{columns}
\end{frame}

\begin{frame}{Branch: uma linha de desenvolvimento}
\begin{columns}[c]
\column{0.6\textwidth}
\begin{tikzpicture}[scale=0.8]
  \node[commit] (a) at (0,0) {};  \node[commit] (b) at (1.3,0) {};
  \node[commit] (c) at (2.6,0) {}; \node[commit] (d) at (3.9,0) {}; \node[commit] (e) at (6.5,0) {};
  \draw[hist] (a)--(b)--(c)--(d)--(e);
  \node[commit=verdeMB] (f1) at (2.6,1.3) {}; \node[commit=verdeMB] (f2) at (3.9,1.3) {}; \node[commit=verdeMB] (f3) at (5.2,1.3) {};
  \draw[hist=verdeMB] (b) to[out=0,in=180] (f1); \draw[hist=verdeMB] (f1)--(f2)--(f3);
  \draw[hist=verdeMB, dashed] (f3) to[out=0,in=180] (e);
  \node[commit=ouroMB!85!black] (g1) at (3.9,-1.3) {}; \node[commit=ouroMB!85!black] (g2) at (5.2,-1.3) {};
  \draw[hist=ouroMB!85!black] (c) to[out=0,in=180] (g1); \draw[hist=ouroMB!85!black] (g1)--(g2);
  \node[rotulo, anchor=west] at ([xshift=2mm]e.east) {main};
  \node[rotulo=verdeMB, anchor=west] at ([xshift=2mm]f3.east) {feature/login};
  \node[rotulo=ouroMB!85!black, anchor=west] at ([xshift=2mm]g2.east) {bugfix/autenticacao};
\end{tikzpicture}
\column{0.36\textwidth}
\begin{itemize}\small
  \item Trabalho \textbf{isolado}, sem mexer na \texttt{main}.
  \item Várias pessoas \textbf{em paralelo}.
  \item Espaço para \textbf{experimentar} e corrigir.
  \item Revisão \textbf{antes} de integrar.
\end{itemize}
\vspace{2mm}
\destaque{Branch \textbf{não} é uma cópia do projeto: é um \alert{ponteiro} para uma linha do histórico.}
\end{columns}
\end{frame}

\begin{frame}{Merge Request: pedir para integrar}
\centering
\begin{tikzpicture}[passo/.style={caixa, minimum width=2.05cm, minimum height=1.25cm, font=\scriptsize}]
  \node[rotulo=verdeMB, font=\small\bfseries] (src) at (0,1.9) {feature/login};
  \node[rotulo, font=\small\bfseries] (dst) at (12.6,1.9) {main};
  \draw[seta, line width=2pt] (src) -- node[above, font=\small\bfseries\color{azulMB}] {Merge Request} (dst);
  \foreach \i/\ic/\t/\cor in {0/\faEye/Visualizar/azulMB, 1/\faCommentDots/Comentar/azulMB, 2/\faRedo/Pedir ajustes/azulMB, 3/\faVial/Testar (CI)/azulMB, 4/\faThumbsUp/Aprovar/azulMB, 5/\faCodeBranch/Merge/verdeMB}{
    \node[passo=\cor] (p\i) at (\i*2.52,0) {{\large\ic}\\[2pt]\t};
  }
  \foreach \i [evaluate=\i as \j using int(\i+1)] in {0,...,4}{\draw[setafina] (p\i) -- (p\j);}
\end{tikzpicture}\\[6mm]
{\footnotesize\textbf{\color{azulMB}Commit}\, registra uma alteração \quad\textcolor{ouroMB}{\faChevronRight}\quad
\textbf{\color{verdeMB}Branch}\, cria uma linha de trabalho \quad\textcolor{ouroMB}{\faChevronRight}\quad
\textbf{\color{azulMB}MR}\, pede a integração}
\end{frame}

\begin{frame}[fragile]{Os comandos do dia a dia}
\footnotesize
\renewcommand{\arraystretch}{1.28}
\begin{columns}[T]
\column{0.48\textwidth}
\begin{tabularx}{\linewidth}{@{}>{\ttfamily\color{azulMB}}l X@{}}
\toprule
\multicolumn{2}{@{}l}{\bfseries\color{verdeMB}Começar e registrar}\\
git clone & copiar um repositório remoto \\
git init & criar um repositório local \\
git status & ver o estado atual \\
git add & preparar alterações \\
git commit & registrar no histórico \\
git log & consultar o histórico \\
git diff & comparar alterações \\
git restore & desfazer mudanças num arquivo \\
\bottomrule
\end{tabularx}
\column{0.48\textwidth}
\begin{tabularx}{\linewidth}{@{}>{\ttfamily\color{azulMB}}l X@{}}
\toprule
\multicolumn{2}{@{}l}{\bfseries\color{verdeMB}Branches e remoto}\\
git branch & listar ou criar branches \\
git switch & trocar de branch \\
git merge & unir históricos \\
git remote & gerenciar remotos \\
git fetch & baixar do remoto, sem mesclar \\
git pull & baixar \textit{e} mesclar \\
git push & enviar para o remoto \\
\bottomrule
\end{tabularx}
\end{columns}
\end{frame}

\begin{frame}{O fluxo básico de trabalho em equipe}
\centering
\begin{tikzpicture}[p/.style={caixa, minimum width=1.5cm, text width=1.3cm, minimum height=1.15cm, font=\scriptsize, inner sep=3pt}]
  \foreach \i/\t/\c in {0/{\ttfamily git\\pull}/azulMB, 1/Alterar\\código/azulMB, 2/{\ttfamily git\\status}/azulMB, 3/{\ttfamily git\\add}/azulMB,
                        4/{\ttfamily git\\commit}/azulMB, 5/{\ttfamily git\\push}/azulMB, 6/Merge\\Request/verdeMB, 7/Revisão\\e merge/verdeMB}{
    \node[p=\c] (n\i) at (\i*1.85,0) {\t};
  }
  \foreach \i [evaluate=\i as \j using int(\i+1)] in {0,...,6}{\draw[setafina] (n\i) -- (n\j);}
  \draw[decorate, decoration={brace, amplitude=6pt, mirror}, azulMedio, thick] (-0.75,-0.8) -- (9.99,-0.8)
    node[midway, below=8pt, font=\scriptsize\color{azulMB}] {no seu computador (Git)};
  \draw[decorate, decoration={brace, amplitude=6pt, mirror}, verdeMB, thick] (10.35,-0.8) -- (13.7,-0.8)
    node[midway, below=8pt, font=\scriptsize\color{verdeMB}] {no GitLab};
\end{tikzpicture}\\[8mm]
\destaque[verdeMB]{\faLightbulb\; Regra de ouro: \textbf{comece sempre com \texttt{git pull}} e trabalhe numa branch própria. Quanto mais tempo longe da \texttt{main}, maior a chance de conflito.}
\end{frame}

\begin{frame}[fragile]{Conflito: quando o Git pede uma decisão humana}
\begin{columns}[T]
\column{0.44\textwidth}
{\small Duas branches mudaram \textbf{a mesma linha}:}\\[2mm]
\begin{tikzpicture}
  \node[caixa=verdeMB, font=\ttfamily\footnotesize] (a) at (0,0) {nome = "João"};
  \node[caixa=ouroMB!85!black, font=\ttfamily\footnotesize] (b) at (3.4,0) {nome = "Maria"};
  \node[font=\scriptsize\bfseries\color{verdeMB}, above=1mm of a] {branch A};
  \node[font=\scriptsize\bfseries\color{ouroMB!80!black}, above=1mm of b] {branch B};
  \node[caixa=vermelhoAlerta, font=\footnotesize] (c) at (1.7,-1.6) {\faQuestionCircle\; qual fica?};
  \draw[setafina] (a) -- (c); \draw[setafina] (b) -- (c);
\end{tikzpicture}\\[3mm]
\destaque{Conflito \textbf{não é erro}. É o Git dizendo: \enquote{não sei combinar isso sozinho, decida você}.}
\column{0.52\textwidth}
{\small O Git marca o arquivo assim:}
\begin{lstlisting}[style=saida]
<<<<<<< HEAD
nome = "João"
=======
nome = "Maria"
>>>>>>> branch-b
\end{lstlisting}
\begin{enumerate}\footnotesize
  \item edite e deixe só o que deve ficar;
  \item apague os marcadores \texttt{<<<}, \texttt{===}, \texttt{>>>};
  \item \cmd{git add arquivo} e \cmd{git commit}.
\end{enumerate}
\end{columns}
\end{frame}

\begin{frame}[fragile]{GitLab CLI: o GitLab pelo terminal}
\begin{columns}[T]
\column{0.5\textwidth}
\begin{lstlisting}[basicstyle=\ttfamily\footnotesize\color{cinzaTexto}]
# autenticar (uma vez)
glab auth login --hostname <endereco-do-gitlab>

# criar um Merge Request da branch atual
glab mr create --target-branch main \
  --title "Integra o sistema"

# acompanhar
glab mr list
glab mr view
glab ci status
\end{lstlisting}
\column{0.44\textwidth}
\begin{tikzpicture}[l/.style={font=\small, anchor=west}]
  \node[l] at (0,1.6) {\textbf{\color{azulMB}Git}\hspace{11mm} controle de versão};
  \node[l] at (0,0.8) {\textbf{\color{azulMB}GitLab}\hspace{5mm} plataforma de colaboração};
  \node[l] at (0,0) {\textbf{\color{azulMB}glab}\hspace{9.5mm} o GitLab no terminal};
\end{tikzpicture}\\[4mm]
\destaque{As \textbf{regras de permissão} do GitLab valem também no terminal: se a \texttt{main} exige aprovação, o \texttt{glab} não pula essa etapa.}
\end{columns}
\end{frame}

% =============================================================================
\secao{02}{Problema 1}{Um sistema já versionado localmente precisa entrar no GitLab institucional, onde a \texttt{main} já existe e está protegida}
% =============================================================================

\begin{frame}{O cenário na instituição}
\begin{columns}[c]
\column{0.5\textwidth}
{\small Ao criar um repositório, o GitLab institucional \textbf{já entrega}:}
\begin{itemize}\small
  \item uma branch \texttt{main} criada;
  \item um \texttt{README.md} padrão;
  \item a \texttt{main} \alert{protegida}: nada de \texttt{push} direto;
  \item mudanças só por \alert{Merge Request}.
\end{itemize}
\vspace{2mm}
{\small Mas o sistema \textbf{já existia} no computador, com anos de histórico Git próprio.}
\column{0.46\textwidth}
\centering
\begin{tikzpicture}
  \node[caixa, minimum width=5.4cm, minimum height=2.2cm, fill=azulClaro] (r) at (0,0) {};
  \node[anchor=north west, font=\scriptsize\bfseries\color{azulMB}] at (r.north west) {\includegraphics[height=3.5mm]{\imagens/gitlab-tanuki.pdf}\; repositório institucional};
  \node[rotulo] (m) at (-1.3,-0.2) {\faLock\; main};
  \node[caixa=ouroMB!85!black, font=\scriptsize\ttfamily] at (1.2,-0.2) {README.md};
  \node[font=\scriptsize\color{vermelhoAlerta}, below=7mm of r] (x) {\faBan\; \texttt{git push origin main} \enspace recusado};
\end{tikzpicture}
\end{columns}
\end{frame}

\begin{frame}[fragile]{O problema: dois históricos que não se conhecem}
\begin{columns}[c]
\column{0.56\textwidth}
\begin{tikzpicture}
  \node[font=\scriptsize\bfseries\color{verdeMB}, anchor=west] at (-0.4,1.05) {REPOSITÓRIO LOCAL (anos de trabalho)};
  \foreach \l/\x in {A/0,B/1.3,C/2.6,D/3.9}{\node[commit=verdeMB] (\l) at (\x,0.35) {\l};}
  \draw[hist=verdeMB] (A)--(B)--(C)--(D);
  \node[font=\scriptsize\bfseries\color{azulMB}, anchor=west] at (-0.4,-0.85) {REPOSITÓRIO GITLAB (recém-criado)};
  \foreach \l/\x in {X/0,Y/1.3}{\node[commit] (\l) at (\x,-1.55) {\l};}
  \draw[hist] (X)--(Y);
  \node[font=\large\color{vermelhoAlerta}] at (2.8,-1.55) {\faUnlink};
  \node[font=\scriptsize, text width=2.6cm, align=left, anchor=west] at (3.3,-1.55) {nenhum ancestral em comum};
\end{tikzpicture}
\column{0.4\textwidth}
{\small Um \texttt{git merge} comum simplesmente recusa:}
\begin{lstlisting}[style=saida]
$ git merge origin/main
fatal: refusing to merge
unrelated histories
\end{lstlisting}
\destaque{O Git só sabe mesclar linhas que \textbf{nasceram de um mesmo ponto}. Aqui, não há esse ponto.}
\end{columns}
\end{frame}

\begin{frame}{O objetivo: um histórico só, sem perder nada}
\centering
\begin{tikzpicture}
  \foreach \l/\x in {A/0,B/1.6,C/3.2,D/4.8}{\node[commit=verdeMB] (\l) at (\x,1) {\l};}
  \draw[hist=verdeMB] (A)--(B)--(C)--(D);
  \foreach \l/\x in {X/0,Y/2.4}{\node[commit] (\l) at (\x,-1) {\l};}
  \draw[hist] (X)--(Y);
  \node[commitm] (M) at (6.8,0) {M};
  \draw[hist=verdeMB] (D) to[out=0,in=135] (M);
  \draw[hist] (Y) to[out=0,in=-135] (M);
  \node[rotulo=verdeMB, anchor=west] at ([xshift=4mm]M.east) {desenvolvimento};
  \node[font=\scriptsize\color{verdeMB}, anchor=east] at ([xshift=-3mm]A.west) {histórico do sistema};
  \node[font=\scriptsize\color{azulMB}, anchor=east] at ([xshift=-3mm]X.west) {histórico institucional};
  \node[font=\scriptsize, align=center, below=4mm of M] {commit de merge\\une os dois};
\end{tikzpicture}\\[6mm]
\destaque[verdeMB]{Depois do commit \textbf{M}, os dois históricos passam a ter ancestral comum. A partir daí um Merge Request para a \texttt{main} flui \textbf{sem conflito}.}
\end{frame}

\begin{frame}[fragile]{Passos 1 a 3: conectar, baixar e renomear}
\begin{columns}[T]
\column{0.55\textwidth}
\begin{lstlisting}
# 1. dar ao Git local um endereço remoto
git remote add origin <URL_DO_REPOSITORIO>
git remote -v

# 2. baixar o que existe no GitLab
#    (sem mesclar nada ainda)
git fetch origin

# 3. tirar a main local do caminho
git branch -m main desenvolvimento
\end{lstlisting}
\column{0.41\textwidth}
\begin{tikzpicture}[l/.style={font=\scriptsize, anchor=west, text width=5cm}]
  \node[rotulo=verdeMB] (d) at (0,0) {desenvolvimento};
  \node[l] at (1.9,0) {seu histórico (A–D), \\ com novo nome};
  \node[rotulo] (o) at (0,-1.2) {origin/main};
  \node[l] at (1.9,-1.2) {cópia de leitura da \\ \texttt{main} institucional};
\end{tikzpicture}\\[5mm]
\destaque{Renomear evita confundir a sua \texttt{main} com a \texttt{main} \textbf{protegida} do GitLab.}
\end{columns}
\end{frame}

\begin{frame}[fragile]{Passo 4: autorizar a união dos históricos}
\begin{lstlisting}
git merge --allow-unrelated-histories origin/main
\end{lstlisting}
\vspace{2mm}
\begin{columns}[T]
\column{0.47\textwidth}
\begin{tikzpicture}\node[caixa=vermelhoAlerta, text width=5.9cm, align=left, font=\small]{%
  {\bfseries\color{vermelhoAlerta}\faTimes\; NÃO significa}\\[2pt]
  \enquote{ignore todos os conflitos}.};\end{tikzpicture}
\column{0.47\textwidth}
\begin{tikzpicture}\node[caixa=verdeMB, text width=5.9cm, align=left, font=\small]{%
  {\bfseries\color{verdeMB}\faCheck\; Significa}\\[2pt]
  \enquote{eu sei que não há ancestral comum e autorizo o Git a tentar juntar}.};\end{tikzpicture}
\end{columns}
\vspace{4mm}
{\small Resultado típico: os dois lados têm um \texttt{README.md} diferente.}
\begin{lstlisting}[style=saida]
CONFLITO (adicionar/adicionar): conflito de mesclagem em README.md
Automatic merge failed; fix conflicts and then commit the result.
\end{lstlisting}
\end{frame}

\begin{frame}[fragile]{Passo 5: \texttt{ours} e \texttt{theirs} — o ponto que mais causa erro}
\begin{columns}[c]
\column{0.5\textwidth}
\centering
\begin{tikzpicture}
  \node[commitm, minimum size=12mm, font=\bfseries\color{azulMB}] (m) at (0,0) {merge};
  \node[caixa=verdeMB, text width=2.7cm] (o) at (-1.9,-2) {\textbf{OURS}\\ \scriptsize a branch em que \textbf{você está}\\ (\texttt{desenvolvimento})};
  \node[caixa, text width=2.7cm] (t) at (1.9,-2) {\textbf{THEIRS}\\ \scriptsize a branch que \textbf{está chegando}\\ (\texttt{origin/main})};
  \draw[setafina] (m) -- (o); \draw[setafina] (m) -- (t);
\end{tikzpicture}
\column{0.46\textwidth}
\begin{lstlisting}
# manter o README do sistema
git restore --ours README.md
# forma antiga equivalente
git checkout --ours README.md

# manter o README institucional
git restore --theirs README.md
\end{lstlisting}
\end{columns}
\vspace{3mm}
\destaque[vermelhoAlerta]{\faExclamationTriangle\; \texttt{ours} \textbf{não} quer dizer \enquote{o arquivo do meu computador}. Quer dizer \textbf{o lado da branch atual naquele merge}. Confunda isso e você preserva o arquivo errado.}
\end{frame}

\begin{frame}[fragile]{Passos 6 e 7: registrar e publicar}
\begin{columns}[T]
\column{0.5\textwidth}
\begin{lstlisting}
# 6. marcar o conflito como resolvido
git add README.md
git commit   # grava o M

# 7. enviar a nova branch ao GitLab
git push origin desenvolvimento
\end{lstlisting}
\column{0.46\textwidth}
{\small No GitLab passam a existir, lado a lado:}\\[3mm]
\begin{tikzpicture}[l/.style={font=\scriptsize, anchor=west}]
  \node[rotulo] at (0,0) {\faLock\; main}; \node[l] at (1.3,0) {intocada e protegida};
  \node[rotulo=verdeMB] at (0.45,-0.8) {desenvolvimento}; \node[l] at (1.9,-0.8) {sistema + histórico integrado};
\end{tikzpicture}
\end{columns}
\end{frame}

\begin{frame}[fragile]{Passo 8: o Merge Request fecha o ciclo}
\begin{columns}[c]
\column{0.5\textwidth}
\centering
\begin{tikzpicture}
  \node[rotulo=verdeMB, font=\small\bfseries] (d) at (0,1.2) {desenvolvimento};
  \node[rotulo, font=\small\bfseries] (m) at (0,-1.2) {\faLock\; main};
  \draw[seta, line width=2pt] (d) -- node[right=2mm, font=\small, align=left] {Merge Request\\[-1pt]\scriptsize revisão + aprovação} (m);
\end{tikzpicture}
\column{0.46\textwidth}
\begin{itemize}\small
  \item Feito na \textbf{interface web} do GitLab, com aprovação de quem tem permissão.
  \item Ou pelo terminal, com o \texttt{glab}:
\end{itemize}
\begin{lstlisting}
glab mr create -s desenvolvimento -b main
\end{lstlisting}
\end{columns}
\vspace{3mm}
\destaque[verdeMB]{Como os históricos agora são relacionados, o merge \textbf{não é recusado}: todos os conflitos já foram tratados localmente.}
\end{frame}

\begin{frame}[fragile]{Resumo do Problema 1 em oito linhas}
\begin{lstlisting}[basicstyle=\ttfamily\footnotesize\color{cinzaTexto}]
git remote add origin <URL_DO_REPOSITORIO>           # 1. conectar
git fetch origin                                      # 2. baixar
git branch -m main desenvolvimento                    # 3. renomear
git merge --allow-unrelated-histories origin/main     # 4. unir
git restore --ours README.md                          # 5. decidir o conflito
git add README.md && git commit                       # 6. registrar
git push origin desenvolvimento                       # 7. publicar
glab mr create -s desenvolvimento -b main             # 8. pedir a integração
\end{lstlisting}
\vspace{2mm}
\destaque{\faBook\; A apostila traz este mesmo caso para você \textbf{reproduzir no seu computador}, sem precisar de servidor.}
\end{frame}

% =============================================================================
\secao{03}{Problema 2}{Os sistemas agora estão versionados. Mas quase nenhum tinha documentação}
% =============================================================================

\begin{frame}{Versionado não é documentado}
\centering
\begin{tikzpicture}[p/.style={caixa, minimum width=2.1cm, text width=1.9cm, minimum height=1.15cm, font=\scriptsize}]
  \foreach \i/\t/\c in {0/Código\\existente/azulMB, 1/Pouca\\documentação/azulMB, 2/Conhecimento\\concentrado/azulMB,
                        3/Manutenção\\difícil/ouroMB!85!black, 4/Evolução\\difícil/ouroMB!85!black, 5/Risco\\operacional/vermelhoAlerta}{
    \node[p=\c] (n\i) at (\i*2.4,0) {\t};
  }
  \foreach \i [evaluate=\i as \j using int(\i+1)] in {0,...,4}{\draw[setafina] (n\i) -- (n\j);}
\end{tikzpicture}\\[7mm]
\begin{columns}[T]
\column{0.46\textwidth}
\begin{itemize}\footnotesize
  \item documentação ausente ou desatualizada;
  \item regras de negócio só na cabeça de alguém;
  \item dependências e configurações desconhecidas.
\end{itemize}
\column{0.46\textwidth}
\destaque{A solução adotada: um \textbf{agente de IA} (Claude Code) para acelerar o entendimento, a refatoração e a documentação profunda.}
\end{columns}
\end{frame}

\begin{frame}{IA: ferramenta de engenharia, não autoridade}
\centering
\begin{tikzpicture}[c/.style={caixa, minimum width=3.4cm, minimum height=2.7cm, align=left, font=\footnotesize, anchor=north}]
  \node[c] (e1) at (0,0) {{\bfseries\color{azulMB}\faUserTie\; Engenheiro}\\[3pt] • define o objetivo\\ • define restrições\\ • orienta o agente};
  \node[c=ouroMB!85!black] (ia) at (4.6,0) {{\bfseries\color{azulMB}\includegraphics[height=3mm]{\imagens/claude-simbolo.pdf}\; Agente de IA}\\[3pt] • analisa\\ • sugere\\ • documenta\\ • executa o autorizado};
  \node[c=verdeMB] (e2) at (9.2,0) {{\bfseries\color{azulMB}\faUserCheck\; Engenheiro}\\[3pt] • revisa\\ • valida no sistema real\\ • decide};
  \draw[seta] (e1.east|-0,-1.35) -- (ia.west|-0,-1.35);
  \draw[seta] (ia.east|-0,-1.35) -- (e2.west|-0,-1.35);
\end{tikzpicture}\\[5mm]
\destaque[vermelhoAlerta]{A IA pode interpretar errado regras de negócio, intenções e comportamentos implícitos. \textbf{Toda documentação gerada por IA é validada contra o sistema real.} A decisão final é sempre humana.}
\end{frame}

\begin{frame}{Como o prompt evoluiu (dados reais)}
\begin{columns}[c]
\column{0.56\textwidth}
\begin{tikzpicture}
\begin{axis}[
  xbar, width=8.4cm, height=5.6cm, bar width=9pt,
  xmode=log, log basis x=10, xmin=100, xmax=250000,
  xlabel={\scriptsize tamanho do prompt (caracteres, escala log)},
  symbolic y coords={V4,V3,V2,V1,V0}, ytick=data, y dir=normal,
  yticklabels={{\bfseries V4 · 15/09},{V3 · 14/09},{V2 · 10/08},{V1 · 16/07},{V0 · 13/07}},
  yticklabel style={font=\scriptsize\color{azulMB}},
  xticklabel style={font=\tiny\color{azulMedio}},
  axis x line*=bottom, axis y line*=left, x axis line style={azulMedio!40}, y axis line style={draw=none},
  xmajorgrids, grid style={azulMedio!15}, major tick style={draw=none},
  nodes near coords, nodes near coords style={font=\scriptsize\color{cinzaTexto}, anchor=west},
  point meta=explicit symbolic, enlarge y limits=0.14]
  \addplot[fill=azulMB, draw=none] coordinates {
    (37883,V4) [37.883] (61556,V3) [61.556] (4551,V2) [4.551] (944,V1) [231–944] (953,V0) [953]};
\end{axis}
\end{tikzpicture}
\column{0.4\textwidth}
\footnotesize
\begin{description}[V0]
\item[\color{azulMB}V0] pipeline em PT: \enquote{não modifique}, descubra antes.
\item[\color{azulMB}V1] pedidos curtos em PT, crescendo a cada reuso.
\item[\color{azulMB}V2] lista estruturada em PT: 47 itens.
\item[\color{azulMB}V3] protocolo em inglês, fase \emph{read-only}.
\item[\color{verdeMB}V4] prompt final em inglês, 42 seções, \textbf{saída obrigatória em PT-BR}.
\end{description}
\end{columns}
\vfill
{\tiny\color{azulMedio} Fonte: histórico de prompts do Claude Code do autor, 10/07 a 28/09/2026 — 717 prompts em 123 sessões. V1 mostra a faixa observada; o 1º pedido livre tinha 231 caracteres.}
\end{frame}

\begin{frame}{Cada falha virou uma regra}
\footnotesize
\renewcommand{\arraystretch}{1.9}
\begin{tabularx}{\linewidth}{@{}>{\raggedright\arraybackslash\leavevmode\color{vermelhoAlerta}}X c >{\raggedright\arraybackslash\leavevmode\color{verdeMB}\bfseries}X@{}}
\toprule
\textbf{\color{azulMB}O que deu errado} & & \textbf{\color{azulMB}A regra que entrou no prompt} \\
\midrule
Documentação superficial, pedida de novo & \faLongArrowAltRight & Análise em fases, arquivo a arquivo \\
Agente preenchia lacunas sozinho & \faLongArrowAltRight & Registrar como \texttt{GAP}; nunca inventar \\
Risco de alterar código na análise & \faLongArrowAltRight & Fase 0 somente leitura (\emph{read-only}) \\
Prompt em inglês, resposta em inglês & \faLongArrowAltRight & Exigência global: tudo em PT-BR \\
Agente seguia em frente diante de risco & \faLongArrowAltRight & Protocolo \emph{Safe Stop-and-Report} \\
\bottomrule
\end{tabularx}
\end{frame}

\begin{frame}{O que melhorou, em números}
\begin{columns}[c]
\column{0.3\textwidth}
\centering
{\titulofonte\fontsize{54}{54}\selectfont\color{vermelhoAlerta}\raisebox{0pt}{\textasciitilde}53\%}\\
{\footnotesize das tarefas de documentação precisaram de \textbf{retrabalho}\\ com prompts livres (V1, V2)\\ \color{azulMedio}8 de 15 tarefas}
\column{0.06\textwidth}
\centering{\Huge\color{ouroMB}\faLongArrowAltRight}
\column{0.3\textwidth}
\centering
{\titulofonte\fontsize{54}{54}\selectfont\color{verdeMB}\textasciitilde25\%}\\
{\footnotesize com prompts estruturados\\ com contrato de execução (V0, V3, V4)\\ \color{azulMedio}1 de 4 tarefas}
\column{0.3\textwidth}
\destaque{\footnotesize\textbf{Leia com cautela:} estimativa com amostra pequena. \enquote{Retrabalho} = refazer, reenviar, \enquote{continue de onde parou}, idioma errado. Parte das paradas veio de limite de uso, não do prompt.}
\end{columns}
\vspace{5mm}
\centering
\begin{tikzpicture}[s/.style={font=\scriptsize, text width=3.8cm, align=center}]
  \node[s] at (0,0) {\textbf{\color{azulMB}164×}\\ crescimento do prompt, do 1º pedido livre ao final};
  \node[s] at (4.4,0) {\textbf{\color{azulMB}5 regras}\\ nasceram de 5 falhas observadas};
  \node[s] at (8.8,0) {\textbf{\color{azulMB}0}\\ pedidos para refazer a doc após o V4 (1 uso)};
\end{tikzpicture}
\end{frame}

\begin{frame}{O método: documentação em fases}
\centering
\begin{tikzpicture}[f/.style={caixa, minimum width=1.85cm, minimum height=1.55cm, font=\scriptsize, text width=1.7cm}]
  \foreach \i/\n/\t/\c in {0/0/Descoberta somente leitura/verdeMB, 1/1/Mapa do sistema/azulMB, 2/2/Análise profunda/azulMB,
                           3/3/Modelo de conhecimento/azulMB, 4/4/Validação do entendimento/ouroMB!85!black, 5/5/Documentação técnica/azulMB, 6/6+/Auditoria e evolução/azulMB}{
    \node[f=\c] (f\i) at (\i*2.1,0) {{\titulofonte\Large FASE \n}\\[1pt]\t};
  }
  \foreach \i [evaluate=\i as \j using int(\i+1)] in {0,...,5}{\draw[setafina] (f\i) -- (f\j);}
\end{tikzpicture}\\[6mm]
\destaque{Evita a cascata: \emph{entendimento errado} \faLongArrowAltRight\ \emph{documentação errada} \faLongArrowAltRight\ \emph{recomendação errada} \faLongArrowAltRight\ \emph{projeto futuro errado}.\\[2pt]
Em vez disso: \textbf{descobrir \faLongArrowAltRight\ evidenciar \faLongArrowAltRight\ modelar \faLongArrowAltRight\ validar \faLongArrowAltRight\ documentar \faLongArrowAltRight\ auditar \faLongArrowAltRight\ evoluir}.}
\end{frame}

\begin{frame}{As três travas de segurança do prompt}
\begin{columns}[T]
\column{0.31\textwidth}
\begin{tikzpicture}\node[caixa=verdeMB, text width=3.9cm, minimum height=5.1cm, align=left, font=\footnotesize, anchor=north]{%
{\bfseries\color{verdeMB}\large\faEye}\\[2pt]{\bfseries READ-ONLY\\DISCOVERY}\\[4pt]
Na descoberta, o agente \textbf{não} altera código, não apaga, não renomeia, não mexe em configuração nem faz commit.\\[4pt]
\textit{Descobrir antes de modificar.}};\end{tikzpicture}
\column{0.31\textwidth}
\begin{tikzpicture}\node[caixa=ouroMB!85!black, text width=3.9cm, minimum height=5.1cm, align=left, font=\footnotesize, anchor=north]{%
{\bfseries\color{ouroMB!80!black}\large\faBalanceScale}\\[2pt]{\bfseries NUNCA\\INVENTAR}\\[4pt]
Todo fato precisa de evidência. O que não se confirma recebe rótulo:\\[2pt]
{\scriptsize\ttfamily DESCONHECIDO · NÃO VERIFICADO · INFERIDO · REQUER VALIDAÇÃO}\\[4pt]
E separa \textbf{AS-IS}, \textbf{TO-BE} e \textbf{GAP}.};\end{tikzpicture}
\column{0.31\textwidth}
\begin{tikzpicture}\node[caixa=vermelhoAlerta, text width=3.9cm, minimum height=5.1cm, align=left, font=\footnotesize, anchor=north]{%
{\bfseries\color{vermelhoAlerta}\large\faHandPaper}\\[2pt]{\bfseries SAFE\\STOP-AND-REPORT}\\[4pt]
Diante de dúvida, contradição, risco ou falta de autorização:\\[2pt]
\textbf{PARAR \faLongArrowAltRight\ REPORTAR \faLongArrowAltRight\ AGUARDAR} a decisão humana.\\[4pt]
Gravidade de STOP-P0 (crítico) a STOP-P3.};\end{tikzpicture}
\end{columns}
\end{frame}

\begin{frame}{Prompt em inglês, documentação em português}
\begin{columns}[c]
\column{0.5\textwidth}
\centering
\begin{tikzpicture}[p/.style={caixa, minimum width=5.2cm, font=\small}]
  \node[p] (a) at (0,0) {Prompt de instrução \textbf{em inglês}};
  \node[p=ouroMB!85!black] (b) at (0,-1.3) {Agente executa a análise};
  \node[p=verdeMB] (c) at (0,-2.6) {Documentação \textbf{em PT-BR}};
  \draw[setafina] (a) -- (b); \draw[setafina] (b) -- (c);
\end{tikzpicture}
\column{0.46\textwidth}
\begin{itemize}\small
  \item Idioma do prompt \(\neq\) idioma do resultado.
  \item O V3, só em inglês, fez o agente \textbf{responder em inglês}.
  \item No dia seguinte, o V4 ganhou a regra:
\end{itemize}
{\footnotesize\ttfamily\color{azulMB} ALL CONTENT YOU PRODUCE MUST BE WRITTEN IN BRAZILIAN PORTUGUESE.}\\[3mm]
{\scriptsize\color{azulMedio} A economia de tokens com o inglês é a motivação adotada; ela não foi medida no histórico.}
\end{columns}
\end{frame}

\begin{frame}{O que o prompt pede e o que ele entrega}
\begin{columns}[T]
\column{0.48\textwidth}
{\small\bfseries\color{azulMB} Seis perguntas que guiam o agente}\\[1mm]
\footnotesize
\begin{tabular}{@{}>{\bfseries\color{verdeMB}}l l@{}}
Descoberta & O que existe neste sistema? \\
Compreensão & Como as partes se relacionam? \\
Funcionamento & Como o sistema funciona? \\
Auditoria & Quais os riscos e pontos fracos? \\
Documentação & Como virar conhecimento útil? \\
Validação & Isso bate com o sistema real? \\
\end{tabular}\\[4mm]
\destaque[verdeMB]{\footnotesize De \enquote{código que poucas pessoas conhecem} para \enquote{sistema que outras pessoas entendem e mantêm}.}
\column{0.48\textwidth}
{\small\bfseries\color{azulMB} Exemplo do que sai em \texttt{docs/}}\\[1mm]
\begin{tikzpicture}[f/.style={font=\scriptsize\ttfamily, anchor=west, inner sep=1.5pt}]
  \node[f, font=\scriptsize\ttfamily\bfseries] at (0,0) {docs/};
  \foreach \t [count=\y] in {{00-discovery/DISCOVERY.md}, {00-discovery/SYSTEM\_MAP.md}, architecture.md, business-rules.md,
                     database.md · api.md, security.md, known-issues.md, technical-debt.md, guia para leigos}{
    \node[f] at (0.35,-\y*0.36) {\color{azulMedio}├─\color{cinzaTexto} \t};
  }
  \node[f] at (0,-3.6) {README.md \color{azulMedio}(raiz, profissional)};
\end{tikzpicture}\\[1mm]
{\scriptsize\color{azulMedio} A estrutura exata se adapta a cada sistema.}
\end{columns}
\end{frame}

% =============================================================================
\secao{04}{Encerramento}{A visão geral e a mensagem final}
% =============================================================================

\begin{frame}{A visão geral}
\centering
\begin{tikzpicture}[p/.style={caixa, minimum width=3.2cm, font=\small}]
  \node[p, fill=azulMB, text=white, font=\small\bfseries] (s) at (0,0) {SISTEMA};
  \node[p=verdeMB] (v) at (-3.6,-1.3) {\textbf{Versionamento}};
  \node[p=ouroMB!85!black] (k) at (3.6,-1.3) {\textbf{Conhecimento}};
  \node[p=verdeMB] (g) at (-3.6,-2.5) {Git + GitLab};
  \node[p=ouroMB!85!black] (i) at (3.6,-2.5) {IA + validação humana};
  \node[p=verdeMB] (gg) at (-3.6,-3.7) {Histórico e revisão};
  \node[p=ouroMB!85!black] (d) at (3.6,-3.7) {Documentação};
  \node[p, fill=verdeClaro, minimum width=5cm] (f) at (0,-4.9) {\textbf{Manutenção e evolução seguras}};
  \draw[setafina] (s) -| (v); \draw[setafina] (s) -| (k);
  \draw[setafina] (v)--(g); \draw[setafina] (g)--(gg); \draw[setafina] (k)--(i); \draw[setafina] (i)--(d);
  \draw[setafina] (gg) |- (f); \draw[setafina] (d) |- (f);
\end{tikzpicture}
\end{frame}

\begin{frame}{Mensagem final}
\begin{columns}[c]
\column{0.55\textwidth}
{\large O conhecimento técnico da equipe \textbf{não pode morar só na cabeça} de quem escreveu o sistema.}\\[5mm]
\begin{itemize}\small
  \item O código tem \alert{histórico}.
  \item O sistema tem \alert{documentação}.
  \item As alterações têm \alert{revisão}.
  \item As decisões têm \alert{contexto}.
\end{itemize}
\column{0.4\textwidth}
\destaque{A IA acelera esse trabalho \textbf{dentro de um processo controlado}: objetivo claro, restrições bem definidas e validação humana.}\\[4mm]
{\titulofonte\Large\color{azulMB} Tornar os sistemas mais fáceis de compreender, manter, compartilhar e evoluir.}
\end{columns}
\end{frame}

{\setbeamertemplate{footline}{}
\begin{frame}[plain]
\begin{tikzpicture}[remember picture, overlay]
  \fill[azulMB] (current page.south west) rectangle (current page.north east);
  \node[text=ouroMB, font=\titulofonte\fontsize{72}{72}\selectfont] at ([yshift=14mm]current page.center) {OBRIGADO!};
  \node[text=white, font=\montlight\Large] at ([yshift=-4mm]current page.center) {Perguntas?};
  \node[text=white!80!azulMB, font=\small, align=center, text width=11cm] at ([yshift=-20mm]current page.center)
    {\faBook\; A \textbf{apostila} traz tudo o que vimos com mais calma, o passo a passo para reproduzir o Problema 1 e os \textbf{exercícios práticos}.};
  \node[anchor=south, inner sep=0pt] at ([yshift=8mm]current page.south)
    {\marca[white]{24}};
\end{tikzpicture}
\end{frame}}

\end{document}
